% ---- pg-545 （书页 533） ----

\[
\begin{aligned}
d_{(135)2}&=\min\{d_{(35)2},d_{12}\}=\min\{7,9\}=7\\
d_{(135)4}&=\min\{d_{(35)4},d_{14}\}=\min\{8,6\}=6
\end{aligned}
\]
我们得到下一个聚类水平的距离矩阵为
\[
\begin{array}{r@{\;}c}
 & \begin{array}{ccc}(135)&2&4\end{array}\\[2pt]
\begin{array}{r}(135)\\2\\4\end{array} &
\left[\begin{array}{ccc}
0&&\\
7&0&\\
6&\tikz{\node[draw,ellipse,inner sep=1pt]{5};}&0
\end{array}\right]
\end{array}
\]
各对聚类之间的最小最近邻近点距离为 $d_{42}=5$, 于是我们将对象 4 和 2 合并, 得到聚类(24).

现在我们有两个不同的聚类(135)和(24). 它们的最邻近点距离(即唯一的距离)为
\[
d_{(135)(24)}=\min\{d_{(135)2},d_{(135)4}\}=\min\{7,6\}=6
\]
最终的距离阵变成
\[
\begin{array}{r@{\;}c}
 & \begin{array}{cc}(135)&(24)\end{array}\\[2pt]
\begin{array}{r}(135)\\(24)\end{array} &
\left[\begin{array}{cc}
0&\\
\tikz{\node[draw,ellipse,inner sep=1pt]{6};}&0
\end{array}\right]
\end{array}
\]
于是将聚类(135)和(24)合并, 形成一个包含全部五个对象的聚类(12345), 此时最近邻近点距离为 6.

图解上述分层聚类过程的树图示于\hyperref[fig:12.3]{图 12.3}. 此图清楚说明了分组与分组发生时的距离水平.\hfill$\blacksquare$

\begin{mvfigure}
\begin{tikzpicture}[x=1.0cm,y=0.42cm,every node/.style={font=\small}]
% 坐标轴
\draw (-0.5,0) -- (-0.5,6.5);
\foreach \ty in {0,1,2,3,4,5,6} { \draw (-0.5,\ty) -- (-0.62,\ty); }
\foreach \ty/\tl in {0/0, 2/2, 4/4, 6/6} { \node[left] at (-0.66,\ty) {\tl}; }
\node[rotate=90] at (-1.30,3) {距离};
\draw (-0.5,0) -- (4.45,0);
% 叶结点
\foreach \lx/\lab in {0/1, 1/3, 2/5, 3/2, 4/4} {
  \node[circle,fill,inner sep=1.4pt] at (\lx,0) {};
  \node[below] at (\lx,-0.30) {\lab};
}
\node at (2,-1.20) {对象};
% 树
\draw (1,0) -- (1,2); \draw (2,0) -- (2,2); \draw (1,2) -- (2,2);
\draw (1.5,2) -- (1.5,3);
\draw (0,0) -- (0,3); \draw (0,3) -- (1.5,3);
\draw (0.75,3) -- (0.75,6);
\draw (3,0) -- (3,5); \draw (4,0) -- (4,5); \draw (3,5) -- (4,5);
\draw (3.5,5) -- (3.5,6);
\draw (0.75,6) -- (3.5,6);
\end{tikzpicture}
\phantomsection\label{fig:12.3}\mvcap{图 12.3\quad 五个对象之间距离的单连接树图}
\end{mvfigure}

在典型的分层聚类应用中, 人们感兴趣的主要是中间结果, 即将对象分为中等数目的聚类.

\noindent\textbf{例 12.4}\quad (11 种语言的单连接聚类)

考虑\hyperref[tab:12.3]{表 12.3} 中的和谐矩阵, 它代表 11 种语言中前 10 个数字词之间的接近程度. 为构造距离矩阵, 我们将和谐性(两种语言的 10 个数字的首字母相同的频数)从完全一致的数 10(每种语言自身的固有频数)中减去, 得到的距离阵为

\begin{mvtable}
\footnotesize
\[
\begin{array}{r@{\;}c}
 & \begin{array}{ccccccccccc}\mathrm{E}&\mathrm{N}&\mathrm{Da}&\mathrm{Du}&\mathrm{G}&\mathrm{Fr}&\mathrm{Sp}&\mathrm{I}&\mathrm{P}&\mathrm{H}&\mathrm{Fi}\end{array}\\[2pt]
\begin{array}{r}\mathrm{E}\\\mathrm{N}\\\mathrm{Da}\\\mathrm{Du}\\\mathrm{G}\\\mathrm{Fr}\\\mathrm{Sp}\\\mathrm{I}\\\mathrm{P}\\\mathrm{H}\\\mathrm{Fi}\end{array} &
\left[\begin{array}{ccccccccccc}
0&&&&&&&&&&\\
2&0&&&&&&&&&\\
2&\tikz{\node[draw,ellipse,inner sep=1pt]{1};}&0&&&&&&&&\\
7&5&6&0&&&&&&&\\
6&4&5&5&0&&&&&&\\
6&6&6&9&7&0&&&&&\\
6&6&5&9&7&2&0&&&&\\
6&6&5&9&7&\tikz{\node[draw,ellipse,inner sep=1pt]{1};}&\tikz{\node[draw,ellipse,inner sep=1pt]{1};}&0&&&\\
7&7&6&10&8&5&3&4&0&&\\
9&8&8&8&9&10&10&10&10&0&\\
9&9&9&9&9&9&9&9&9&8&0
\end{array}\right]
\end{array}
\]
\end{mvtable}
